\begin{scanpage}{0019}{2}
paper of Oberhettinger [20]. Since the problem of diffraction by a wedge lies at the base of our studies, we considered it advisable not only to present the results of its rigorous solution, but also to give them a new more graphic derivation. The idea for this derivation follows directly from the work of Sommerfeld. Sommerfeld found the solution to the problem in the form of a contour integral, and then he transformed it to a series. However, one may proceed in the opposite direction: first find the solution in the form of a series and then give its integral representation. Such a path seems to us more graphic, and is discussed in this section. The necessity for a detailed derivation is caused by the fact that the results of Sommerfeld [16] are not represented in a sufficiently clear form, which hinders their use.

Let us assume there is in free space (a vacuum) an ideally conducting wedge and a cylindrical wave source $Q$ parallel to its edge (Figure 1). Let us introduce the cylindrical coordinate system $r,\phi,z$ in such a way that the $z$ axis coincides with the wedge edge, and the angle $\phi$ is measured from the irradiated surface. The external wedge angle will be designated by the letter $\alpha$, so that $0\leq\phi\leq\alpha$. The coordinates of the source $Q$ we will designate by $r_0,\phi_0$.

Let us investigate two particular cases for the excitation of an electromagnetic field. In the first case, it is excited by a ``filament of electric current''
\[j_z^e=-i\omega p_z\delta(r-r_0,\phi-\phi_0),\tag{1.01}\]
in the second case, it is excited by a ``filament of magnetic current''
\[j_z^m=-i\omega m_z\delta(r-r_0,\phi-\phi_0).\tag{1.02}\]
The quantities $p_z$ and $m_z$ here designate, respectively, the electric and magnetic moments of the filament per unit length along the $z$ axis, $\omega$ is the cyclic frequency $\left(\omega=ck=c\frac{2\pi}{\lambda}\right)$, $\delta(r-r_0,\phi-\phi_0)=\delta(r-r_0)\delta[r(\phi-\phi_0)]$ is a two-dimensional delta function which satisfies the condition
\end{scanpage}
